\subsection{A cache of reasoning, kept as code}
\label{sec:mem:arc}
The memories above give a model something to read; none spares it the
thinking. A KV prefix is reused only up to the first token that differs, and
a stored answer to ``how many divisors does 196 have'' is wrong for 144. After
CPU-ARC \cite{belcak2026cpuarc}, the reasoning cache keeps what System~2
worked out as \emph{code} whose inputs are the question's values, so that a
later question that is the same problem in other words or with other numbers
is answered by running the code. The roles of the paper map onto the pair:
System~2 is both the source and the builder, System~1 reads values out of a
question, and the exact verifiers take the place of judge models.

\emph{Entries.} An entry holds the question it was built from, a template
with one placeholder per input (\texttt{\{n\}}), the inputs (name, type,
description, original value) and \emph{semantic functions}: Python functions
that read their inputs with \texttt{load(ctx, "<description>", type)}, call
one another with \texttt{do(ctx, "<what to do>", payloads)}, record each step
with \texttt{narrate}, and return a dictionary with natural-language keys.
Names never connect functions; descriptions do. Unlike CPU-ARC's, the
functions call no model: every value a step needs is read before the code
runs, so the code is deterministic and runs in the verifier sandbox with no
network. The embedder links each \texttt{do()} description to a function when
the entry is registered; at run time descriptions match exactly, then through
that table, then by word overlap of at least 0.5.

\emph{Building.} When System~2 has answered a maths question with an answer
that was checked or agreed, the question joins a queue that is worked through
only while no request is in flight (a request that arrives stops the build and
requeues it). System~2 receives its own question, the end of its reasoning,
its answer and the four existing function descriptions closest to the
question, and writes three sections: the template, the inputs and the code.
Each failed check is fed back for up to three rounds: the template filled
with the values must be the question (when only the words differ, the template
is rebuilt from the question's own words by replacing each value written
exactly once); the code must pass a static check (standard-library imports
only, no file, \texttt{eval} or \texttt{print}), have exactly one root
function and narrate in every function; the root on the original values must
give the expected answer; the code must still run when each numeric input is
changed (it computes, it does not recall); and the question must find the new
entry with its own values.

\emph{Answering.} The cache is tried before the System~1 draft. An entry is
found in three ways, in order: its template as a pattern, whose captures are the
values; the question's shape with its numbers masked, the numbers then taken
by position; and the meaning vector of the question. The first two match only the whole of
a stored, self-contained question, so they run for any question; the third is
kept to maths questions that do not lean on the previous turn (after MeanCache
\cite{gill2025meancache}). On the meaning tier two checks come first, because the vector finds the same topic, not only the same
problem: a short list of word pairs that turn one question into another
(area and perimeter, odd and even, simple and compound interest) refuses the
match without a model, and System~1 is asked in one word whether the new
question asks for the same quantity as the entry's; then it reads the values
as JSON. Two guards apply to every tier: each value must be written in the
question, and each number in the question must be used by some value, so a
question that adds a 2-dollar eraser is not the cached problem. The code then
runs in the sandbox; any error or failed conversion is a miss, and a miss goes
to the models as before.

\emph{Trust.} A new entry is a \emph{candidate}: it runs beside the models,
which still answer, and its result is compared with theirs (with a reference
answer when the request has one; an unchecked quick answer is no evidence).
After two agreements and no disagreement it is \emph{trusted} and answers on
its own, labelled as cross-checked, or verified when the request's reference
confirms it. A disagreement on a checked answer found by the template or the
shape stops the entry. Each entry also learns its own similarity floor for the
meaning tier, lowered to the lowest similarity at which it agreed and kept
above the highest at which it disagreed (a simplified vCache
\cite{schroeder2026vcache}; trust before serving follows Krites
\cite{singh2026krites}). Past 500 entries the least used one that no other
entry calls is removed \cite{kulkarni2026eviction}.

\emph{Compilation.} A mature cache is written into one standard-library
module (CPU-ARC's static resolution): \texttt{load(ctx, d)} becomes a named
parameter, \texttt{do()} a direct call with the inputs a callee was not given
forwarded from its caller, and \texttt{load(result, d)} a key of the callee's
returned dictionary. The module carries each entry's template, so it answers a
question that matches one with no model at all, and it checks itself: an entry
whose compiled form does not give its expected answer is left out, with the
reason.
